\begin{statement}{P3K.33}{Unresolved; scalar-fan audit correction}
\par\noindent\textbf{Audit correction, 17 November 2026.} The explicit parameter/root/cap descent and individual reference fixed-star comparison are retained. The all-fan application of P3K.31 and its numerical range are withdrawn; a simultaneous equivariant product of the full marked plumbing cannot be supplied merely by the fixed-star product. The formulas and earlier argument below are retained as the audited prior claim, not as current proved conclusions. See P3K.61--62 and the correction ledger.
\medskip

Retain the actual K11 construction with logarithmic orders \(3K,4L\),
\(\gcd(3K,4L)=1\), and either sphere datum
\[
       4Ln_0+3Kn_1=\varepsilon,\qquad \varepsilon\in\{1,-1\}.
\]
The additional delta coefficients of the finite logarithmic vectors may
be arbitrary among the already admissible data; they are fixed by the
parameter shear and do not enter the calculation below. Put
\[
 P=4Ln_0,\qquad Q=3Kn_1,\qquad b=2KL,\qquad d=12KL.
\]
The actual closed parameter is
\[
 w_h=e^{-\pi i c/(KL)},\qquad u=\kappa w_h^{2KL},
 \qquad w_h=r^6,\qquad v=\frac{\log r}{2\pi i}.
\]
Thus a positive \(w_h\)-loop sends \(c\) to \(c-b\) and
\(v\) to \(v+1/6\).

There is an explicit extending endpoint translation and an actual
simultaneously closed Cox product on all three cusp curves and both
triple points for this single loop. In that product the fixed-edge
normal integers are
\[
          (k,\ell)=(-KL,-4KL),\qquad
          (m_{\rm rel},n_{\rm rel},s)=(KL,0,0).
\]
For every allowed index-one unimodular periodic fan \(T(v_1,v_2)\),
write the fixed physical circle as \(a v_1+b_1v_2\), with
\(\gcd(a,b_1)=1\), and set
\[
                 S=a^2-ab_1+b_1^2.
\]
The actual positive single-loop return has oriented class
\[
       [\Sigma_{\Phi,T}]
          =KL\frac{(S-1)(S-3)}8\pmod {28}.                 \tag{P3K.33.1}
\]
For the original \(K=L=1\) family the exact set of classes is
\[
 \{0,1,3,6,7,8,10,13,14,15,17,20,21,22,24,27\}.
\]
The same values apply to both determinant signs, by the actual signed
page and Thom maps below. No homotopy-sphere class is divided in the proof.
\end{statement}
\begin{proof}[Preserved earlier argument; scope corrected above]
\par\medskip\noindent\textbf{1. The actual integral one-period primitive.}\enspace

Use exactly the homology marking of P5.6 and K11:
\[
\begin{split}
 A_1&=\begin{pmatrix}
 1&0&0&0\\6&0&1&0\\-6&-1&-1&0\\-2&1&0&1
 \end{pmatrix},\\
 A_2&=\begin{pmatrix}
 1&0&0&0\\0&0&-1&0\\-6&1&0&0\\3&0&1&1
 \end{pmatrix},\qquad
 A_0=(A_1A_2)^{-1}
 =\begin{pmatrix}
 1&0&0&0\\0&1&0&0\\0&1&1&0\\-1&0&0&1
 \end{pmatrix}.
\end{split}                                                     \tag{P3K.33.2}
\]
Here \(e_1\) is the marked gamma period and \(\delta=e_4\).
Let \(J\) have columns \(e_3,\delta\), and put
\(N_1=1+A_1+A_1^2\), \(N_2=1+A_2+A_2^2+A_2^3\).
For clarity, retain the full degree-zero and degree-one maps,
including the finite cyclic group arrows and the cusp invariant stalk:
\[
\begin{split}
 d^0(a,b_0,b_1,z)={}&((A_1-I)a,(A_2-I)a,
 (A_1-I)b_0,(A_2-I)b_1,a-b_0,a-b_1,a-Jz),\\
 d^1(p,q,r,s,u_0,u_1,w)={}&(KN_1r,LN_2s,
 p-r-(A_1-I)u_0,q-s-(A_2-I)u_1,\\
 &\hspace{30mm}-A_0p-A_0A_1q-(A_0-I)w).
\end{split}                                                     \tag{P3K.33.3}
\]
These are K11's actual root-disc and annulus restriction maps.
The first two target summands record the finite-root characteristic
classes, and the last three record the literal annular matches.

Define
\[
\begin{split}
 p=r&=\frac{P\delta-\varepsilon(A_1-I)e_1}{6}
       =(0,-\varepsilon,\varepsilon,(P+2\varepsilon)/6)^t,\\
 q=s&=\frac{Q\delta-\varepsilon(A_2-I)e_1}{6}
       =(0,0,\varepsilon,(Q-3\varepsilon)/6)^t,\\
 u_0&=u_1=w=0,
 \qquad x=(p,q,p,q,0,0,0).
\end{split}                                                     \tag{P3K.33.4}
\]
This is an integral cochain, with zero cusp-annulus component.
Indeed \(P+Q=\varepsilon\), \(Q\) is divisible by three, and \(P\)
is divisible by four. Reducing the sphere equation modulo three
gives \(P\equiv\varepsilon\pmod3\), hence
\(P+2\varepsilon\equiv0\pmod6\). The odd sphere equation
gives \(Q-3\varepsilon\equiv0\pmod6\).

The finite norm equations are literal:
\[
 N_1p=\frac P2\delta,\qquad
 N_2q=\frac{2Q}3\delta.
\]
Moreover
\((A_1A_2-I)e_1=\delta\), so
\[
 p+A_1q=\frac{(P+Q)\delta-\varepsilon(A_1A_2-I)e_1}{6}=0.
\]
Consequently
\[
 d^1x=(bn_0\delta,bn_1\delta,0,0,0),                 \tag{P3K.33.5}
\]
exactly the characteristic change of the actual \(c\mapsto c-b\)
shear. This proves the full local-to-global integral match; it
does not replace a finite norm by an invariant projector.
Every block in \(x\) has first coordinate zero. Hence the actual
successive parameter shears \(S_b=I+bE_{41}\) fix \(x\).

The sixfold endpoint used in K11 has the logarithmic primitive
\[
 x_{\rm cov}=(P\delta,Q\delta,P\delta,Q\delta,0,0,\varepsilon e_1).
\]
There is the following exact integral comparison:
\[
 x_{\rm cov}-6x=d^0y,\qquad
 y=(\varepsilon e_1,\varepsilon e_1,\varepsilon e_1,
       0\in\mathbb Z^2).                             \tag{P3K.33.6}
\]
In particular the cusp component of \(y\) is zero. Formula
(P3K.33.6) follows by substituting (P3K.33.4) in every block; its last
block is \(\varepsilon e_1\), exactly the cusp component of
\(x_{\rm cov}\).

\par\medskip\noindent\textbf{2. Chartwise holomorphic endpoint maps.}\enspace

This cochain is realized by an actual translation on the family.
Use the same integral Jacobian identification
\(\Pi_c=\Pi_{c-b}S_b\) as K11. On the ordinary base, let
\[
 H(t)=-\frac{P\log t+Q\log(t-1)}{2\pi i}.
\]
The cover's ordinary endpoint primitive is \(H\delta\).
At the cusp, with \(q=1/t\) and
\(s=\log q/(2\pi i)\), K11 gives
\[
 H=\varepsilon s+U(q),\qquad
 \gamma_c=(\Gamma(q),C(q)-s-dv)^t,
\]
and its actual bounded cusp primitive is
\(H\delta+\varepsilon\gamma_c\).
All logarithms and the marked period \(\Gamma\) are those of
the original family; a sixth root of a different source
translation is not substituted.

Let \(\tau_{{\rm cov},i}\) denote those actual cover primitives
on the ordinary, two finite-root, and cusp charts. The finite-root
primitives are the holomorphic vectors in K11's extending cover
map. Define the one-step primitives by the explicit formula
\[
       \tau_i=\frac{\tau_{{\rm cov},i}+\Pi_cy_i}{6}.
                                                               \tag{P3K.33.7}
\]
On a finite-root disc all periods in \(\Pi_c y_i\) are
holomorphic in the actual smooth root model, so (P3K.33.7) is
holomorphic there. The change of the annular gluing under the
parameter shift is linear in \(b\) and is in the fixed delta
direction. Thus its sixfold value is exactly six times its
one-step value. Taking the lattice-valued coboundary of (P3K.33.7)
and using (P3K.33.6) gives the integral cochain \(x\), rather
than a fractional cochain. This verifies equivariance for every
finite group arrow as well as equality on every annular overlap.
Exponentiation therefore gives an actual holomorphic isomorphism
\(D_c:X_c\longrightarrow X_{c-b}\). Its inverse is given
chartwise by the negative translation in the target marking.

One can see the finite-disc cancellation directly. The coefficients
of the potentially unbounded logarithms in \(H/6\) are
\(-P/6=-bn_0/(3K)\) and
\(-Q/6=-bn_1/(4L)\). The change of the respective logarithmic
gluing adds their negatives. The remaining unit logarithm and
root-model period in (P3K.33.7) are holomorphic. The finite cyclic
norms in (P3K.33.5) are the exact integral compatibility at the
centres of those discs.

The cusp formula is especially simple because \(y_{\rm cusp}=0\):
\[
 \tau_{\rm cusp}
     =\frac{H\delta+\varepsilon\gamma_c}{6}
     =\binom{\varepsilon\Gamma/6}
              {(\varepsilon C+U-\varepsilon dv)/6}.    \tag{P3K.33.8}
\]
It is a bounded holomorphic vector. In the actual cusp torus
characters it is literally
\[
 (Z,V)\longmapsto
 \left(e^{2\pi i\varepsilon\Gamma(q)/6}Z,
       e^{2\pi i(\varepsilon C(q)+U(q)-\varepsilon dv)/6}V\right).
                                                               \tag{P3K.33.9}
\]
Every Cox monomial \(q^\alpha Z^\beta V^\gamma\) is multiplied
by the corresponding character of these two units. This formula
extends simultaneously to every cone of the chosen fan, fixes
its ray and curve labels, and has the same expression on the
common Cox refinement. It therefore includes all three curves
and both triple points, not just the fixed sphere.

There is also an exact norm check. Under the next parameter
step the marked gamma period changes by \(-b\delta\), so
\[
 \tau(c-b)-\tau(c)=-\frac{\varepsilon b}{6}\delta.
\]
All other cochain blocks remain fixed by \(S_b\). Hence
\[
 \sum_{j=0}^5\tau(c-jb)=6\tau(c)-\frac{5\varepsilon b}{2}\delta
                      =6\tau(c)-5\varepsilon KL\delta. \tag{P3K.33.10}
\]
The last summand is an actual integral delta period. Together
with (P3K.33.6), this shows that \(D^6\) is the previously specified
cover endpoint map in the torus quotient.

The construction also fixes any possible endpoint ambiguity.
For the higher root base, \(H^1(\widehat B,\mathcal H)=0\)
by the same integral calculation as P5.5: multiplication of the
finite norm maps by \(K,L\) does not change their kernels in
the free lattice, and all annular and cusp maps are unchanged.
The source-period principal-cocycle matrix has determinant
\(-1\), as displayed there. Thus an extending translation has
a global Lie logarithm, and its only period is \(\mathbb Z\delta\).
The resulting translation group is \(\mathbb C/\mathbb Z\),
the global physical-delta flow. Therefore changing an admissible
endpoint primitive changes it by an actual extending delta
isotopy. The return diffeomorphism class and the free normal
comparison below are unaffected.

\par\medskip\noindent\textbf{3. A closed original-loop product on every cusp chart.}\enspace

K11's bounded cusp origin on the universal parameter cover is
\[
 \sigma(v)=
 \binom{\varepsilon\Gamma v}
       {v(\varepsilon C+U)-\varepsilon dv^2/2+\varepsilon dv/2}.
\]
Direct subtraction, retaining the constant term, gives
\[
 \sigma(v+1/6)-\sigma(v)
      =\tau_{\rm cusp}+\frac{5\varepsilon KL}{6}\delta. \tag{P3K.33.11}
\]
Define instead the explicit origin
\[
             \sigma_{\rm desc}(v)=\sigma(v)+\varepsilon KLv\delta.
                                                               \tag{P3K.33.12}
\]
Its one-step difference is
\[
 \sigma_{\rm desc}(v+1/6)-\sigma_{\rm desc}(v)
                    =\tau_{\rm cusp}+\varepsilon KL\delta. \tag{P3K.33.13}
\]
In coordinates translated by \(\sigma_{\rm desc}\), the actual
endpoint translation (P3K.33.8) is therefore \(-\varepsilon KL\delta\).
It is literally the identity on every integral torus character.
Thus (P3K.33.12) is a closed simultaneous Cox product for the single
positive \(w_h\)-loop. No residual source translation, curve
permutation, base power, or unrecorded constant remains.

At intermediate parameter angles the two origins differ by the
global physical-delta action. On the sixfold cover the difference
is the closed loop \(\varepsilon KLv\delta\). Its action on the actual K23
quaternion normal coordinate is
\[
             q_{\rm H}\longmapsto
                       q_{\rm H}e^{-2\pi i\varepsilon KLv}.
\]
Consequently the Hopf normal vector
\(q_{\rm H}i\overline{q_{\rm H}}\) is unchanged at every
point of the fixed sphere and every parameter angle. The same
statement on the other two curves is the literal Cox-character
formula (P3K.33.9); those charts are already closed by (P3K.33.13).
This proves the required simultaneous normal endpoint comparison.

The remaining bounded torus units have the exact normal action
computed in K27: in a first chart \(q=\zeta c_0c_1\), a unitary
map \((\zeta,c_0,c_1)\mapsto(a\zeta,pc_0,q_1c_1)\) with
\(apq_1=1\) acts on \(q_{\rm H}\) by
\(q_{\rm H}\mapsto A q_{\rm H}B\), where
\(A^2=pq_1\), \(B^2=p/q_1\). These factors are independent
of the fixed-curve point. Their positive parts are applied to
both compared products. The induced Hopf-normal map is a
constant rotation; its possible endpoint loop extends over the
fixed three-disc and has zero free framing component.

On the ordinary annulus the known K11 difference of origins
has gamma part \(\varepsilon\theta\gamma\). The change (P3K.33.12) is pure
physical delta and is the identity on the orbit. Thus the
smooth node comparison is still exactly K27's
\(R_{\varepsilon\theta}\Phi\), with its positive triangular normal
derivative. This retains the extension of the annular
\(|q|^\theta\) factor through both node charts. On the
sixfold cover it is the signed rotating comparison typed in the
next section.

The closed single-loop Cox framing and the descended signed rotating
framing therefore have the same sixfold free normal class.
Their difference is an integer in the free summand of
\(\pi_1\operatorname{Map}(S^2,SO(3))\); sixfold pullback
multiplies that integer by six, which is injective. The possible
constant rotation is the torsion summand already shown to extend
over the fixed disc. Hence the free classes agree on the single
loop itself. The signed pair proved in the next section is therefore
\[
                       (k,\ell)=(-KL,-4KL)
\]
in the actual simultaneous Cox product (P3K.33.12).

\par\medskip\noindent\textbf{4. The signed original-family dictionary and its normal pair.}\enspace

We identify the original family with K28's actual E6 row, rather
than reversing a different sphere datum. Use the literal holomorphic
base coordinate
\[
             T=\frac{t-1}{t}.
\]
It sends the original \(IV^*\) fibre at \(t=0\) to infinity,
the original III fibre at \(t=1\) to zero, and the selected nodal
fibre at \(t=\infty\) to \(T=1\). The marked data become
\[
\begin{array}{c|c|c}
 &\text{infinity in }T&\text{zero in }T\\\hline
 \text{root order}&m_\infty=3K&m_0=4L\\
 \text{norm coefficient}&n_\infty=n_0&n_0^{\rm new}=n_1
\end{array}
\]
Thus \(m_0n_\infty+m_\infty n_0^{\rm new}=\varepsilon\)
and the actual page degree is \(M=m_\infty m_0=12KL\).
On each root chart this coordinate map multiplies the old base
coordinate by a holomorphic unit. Taking its specified local root
gives the exact root-chart map; it commutes with the cyclic cut.
It has no integral angular winding, so retains the marked norm
coefficients and all finite affine line characters.

The positive endpoint multiplier has exactly the K28 arrow. Indeed
\[
 H=-Q\frac{\log T}{2\pi i}
       +\varepsilon\frac{\log(1-T)}{2\pi i}.
                                                               \tag{P3K.33.14}
\]
K28's positive deck is
\(T^{-Q}(T-1)^{\varepsilon}\), and (P3K.33.14) differs only by
the specified constant factor \((-1)^{\varepsilon}\). Its chosen
logarithm is carried in the common bounded unit of both origins.
At the cusp \(q_T=T-1=-q\); the corresponding constant logarithm
changes \(C\) by the same bounded constant. Therefore K30's
physical coordinates are literally obtained by
\[
 g=\Gamma,\qquad -C_{30}=C+\frac{\log(q_T/q)}{2\pi i},
 \qquad \lambda=u=\kappa r^M.
\]
Its outward gamma is the same marked period \(\gamma_c\).
The line and translation data are exactly the E6 row
\((a_{\rm line},b_{\rm source})=(1,6)\): the source line is
\(\mathcal O(G-O)\), the contracting line map covers
translation by \(-6G\), and its inverse outward gamma period
covers translation by \(+6G\),
and the two good roots are of orders \(3K,4L\). The strict
section/root line maps used for this row in K28 are the same
ones used in K11 and the original source-page construction.

To retain the sign, write \(\alpha=\arg T/(2\pi)\). The actual
solid-torus maps of K28 have ordinary phases
\[
 \phi_1=\varepsilon(m_0\gamma-n_0^{\rm new}\alpha),\qquad
 \phi_2=\varepsilon(m_\infty\gamma+n_\infty\alpha).
                                                               \tag{P3K.33.15}
\]
Their determinant is \(\varepsilon\), and
\(m_0\phi_2-m_\infty\phi_1=\alpha\). They extend on the
actual norm cut \((z,u_i)\mapsto(\zeta_{m_i}z,u_i-n_i)\)
by K28's explicit root-solid-torus formulas. In particular the
positive weighted Hopf angle is \(\beta=\varepsilon\gamma\).
The positive \(r\)-loop has exactly that angular motion by
(P3K.33.14) and K11's displayed origins. Its complex page is the
actual sheet graph
\[
                 \eta^{m_0}-\tau\xi^{m_\infty}=1,
\]
with precisely the root source and line, because the cut changes
the sheet index by \(-\varepsilon m_jn_i\). The sheet-zero
maps of K28 are independent of this sign and retain every
finite character. Thus this is the same marked complex page
and source line for either determinant sign.

At its binding, if \(w\) is the complex parameter at the single
point at infinity, the actual boundary angle is
\(-\arg w=\beta=\theta\), where \(\theta\) is the positive
physical parameter angle. Consequently
\(\partial V=+E\times S^1_\theta\). Orient the actual compression
by \(\partial W=F-E\); the glued trace then has
\(\partial X=+F\times S^1_\theta\). In its actual Hopf chart
the ambient and trace orientations are, respectively,
\[
 \operatorname{or}(F)\operatorname{or}(\nu_{\rm Hopf})\,dr\,d\theta,
 \qquad \operatorname{or}(F)\,dr\,d\theta.
\]
Passing the three normal directions past the last two gives a
positive sign. Hence its meridian sphere is the positive Hopf
sphere and \(\pi_*x=+1\) for both signs. This is the actual
Thom comparison, including the orientation change in (P3K.33.15).

The initial Cox/rotating normal map also retains that sign.
In the beta product its raw matrix is
\(T_\varepsilon Q_x\), where
\(T_\varepsilon=\operatorname{diag}(1,1,\varepsilon)\) and
\(Q_x\) is K30's positive triangular node derivative. The same
\(T_\varepsilon\) is applied to both initial frames. Returning
to their common reference leaves \(Q_x\); conjugation by
\(T_\varepsilon\) is also conjugation by
\(\varepsilon T_\varepsilon\in SO(3)\). Thus neither a hidden
orientation reversal nor a free framing term remains.

For this actual page the good roots contribute no singular Euler
term. The source Euler numbers are \(8,3,1\), so
\[
       \chi=\frac{M(12-8-3)}{12}=KL,
       \qquad H_{\rm page}=M h(G)=M/6=2KL.
\]
These are the line and tangent numbers already proved in K11.
The full degree-\(M\) normal clutch of K30 applies through
the explicit maps (P3K.33.14)--(P3K.33.15). Its two factors are pulled
back from \(T^2\) and \(S^1\), so its cubic clutch integral
is zero. The positive-meridian sphere calculation gives
\[
              k_{\rm cov}=-3H_{\rm page}=-6KL,
              \qquad \ell_{\rm cov}=-24\chi=-24KL,
                                                               \tag{P3K.33.16}
\]
for both \(\varepsilon\). The preceding simultaneous closed
single-loop product proves that these are the pullbacks of its
individual integers. Characteristic-number naturality for the
degree-six parameter cover gives
\(6k=-6KL\), \(6\ell=-24KL\). These are equalities of
ordinary integers in the explicitly compared normal product;
they do not divide a class in \(\Theta_7\). Thus the single-loop
pair is \((-KL,-4KL)\) for either sphere sign. Its integer
index and its initial oriented central class are zero.

\par\medskip\noindent\textbf{5. Applying the proved common-exterior recurrence.}\enspace

The initial vertical-edge return has integer index zero by (P3K.33.16),
so its oriented central class is zero by the already proved Hopf
transfer. K23's exact normal dictionary gives
\[
 m_{\rm rel}=-\ell/4=KL,\qquad
 n_{\rm rel}=-(4k-\ell)/4=0.
\]
The actual hypotheses used in K31 are now all supplied for this
single loop: the common exterior is fixed, all three local Cox
products close simultaneously, their quaternion frame is the
specified one, and the base loop is the original positive
\(w_h\)-loop. K31's explicitly proved recurrence therefore
starts from \((m,n,s)=(KL,0,0)\) and gives
\[
 s(S)=-\frac{KL}{2}(S-1)+\frac{KL}{8}(S^2-1)
      =KL\frac{(S-1)(S-3)}8.
\]
This proves (P3K.33.1) directly on the single loop. In particular the
original family has value zero at \(S=1,3\) and value three
at \(S=7\). An explicit such fan in the original cocharacter
lattice is
\[
 v_1=(3,1),\qquad v_2=(-1,0),\qquad
 \det(v_1,v_2)=1,\qquad
 \delta=(0,1)=v_1+3v_2.
\]
Here \(S=1-3+9=7\), so the original loop has class \(3\)
and Eells--Kuiper invariant \(3/28\). The initial vertical-edge
integer equations are explicitly
\[
 6k=-6,\qquad6\ell=-24,\qquad
 A=\frac{4(-1)-(-4)}{24}=0;
\]
the class \(3\) follows from the actual fan recurrence starting
with those single-loop integers. Its sixfold cover has the corresponding value
eighteen at \(S=7\), but that numerical divisibility is not used
to assign the single-loop class.

The primitive fan values of \(S\) are exactly the odd residue
classes modulo \(112\), by K31's integral basis and Chinese
remainder argument. Evaluating the displayed integer expression
on those fifty-six residues gives the stated sixteen-element set
when \(KL=1\), and multiplication by \(KL\) gives the higher
set. The executable certificate records the entire symbolic
complex, every equality (P3K.33.4)--(P3K.33.13), eight higher integral
tuples for each sphere sign, and this exact residue enumeration in
\path{code/part3_fourfold/original_cox_descent_universal.py},
\path{code/part3_fourfold/original_cox_descent_universal.json}, and
\path{code/part3_fourfold/original_cox_descent_universal.out}.
\end{proof}
